\documentclass[a4paper,12pt]{report}
\usepackage[utf8]{inputenc} 
\usepackage[T1]{fontenc}   
\usepackage[english,french]{babel} 
\usepackage{lipsum}        
\usepackage{graphicx}
%Sommaire clickable: 
\usepackage{color}  
\usepackage{hyperref}
\hypersetup{
    colorlinks=true, 
    linktoc=all,     
    linkcolor=black,  
    urlcolor=blue,
}

\begin{document}

\title{
\includegraphics[width=0.2\textwidth]{CYTech.jpg} \\ Rapport Projet INFO3 P2}
\author{Serrano Lucas Zak Hill}
\date{2024--2025}
\maketitle

\tableofcontents

\chapter*{\vspace*{-2em} Introduction}

\addcontentsline{toc}{chapter}{Introduction}
Ce projet est le projet d'INFO3 de semestre 1 de PRE-ING2 de Cytech (ex EISTI) \\ \\
Le but de ce projet est de réaliser \textbf{une calculatrice} simplifié en C dans le terminal sous Linux, supportant la notation préfixe et postfixe, permettant de réaliser addition, soustraction et multiplication sur des nombres entiers. \\ \\
Se projet est réalisé par Serrano Lucas et Hill Zak, étudiant en PRE-ING2 à Cytech (campus de Pau) du groupe \textit{Cyapp dev}

\chapter*{\vspace*{-2em} Répartition des tâches du projet}
La répartition des tâches du projet a était fait dans l'optique de diviser le travail à réaliser de manière équitable. Voici le tableau ci-dessous détaillant le développeur responsable de chaque grosse partie du projet: \\

\begin{center}
\begin{tabular}{|c|c|}
\hline
\textbf{Tâche} & \textbf{Développeur} \\ \hline
Notation préfixe & Zak Hill \\ \hline
Main & Zak Hill \\ \hline
Mode intéractif & Zak Hill \\ \hline
Arbre & Zak Hill \\ \hline
HelpPage (en roff) & Serrano Lucas \\ \hline
Notation postfixe & Serrano Lucas \\ \hline
Liste & Serrano Lucas \\ \hline
Rapport & Serrano Lucas \\ \hline
\end{tabular}
\end{center}


\chapter*{\vspace*{-2em} Planning}
\addcontentsline{toc}{chapter}{Planning}

Voici un apercu du planning du projet 

\begin{center}
\begin{tabular}{|c|c|}
\hline
\textbf{Semaine} & \textbf{Tâche à faire} \\ \hline
Semaine 1 & Découverte du projet \\ \hline
Semaine 2 & Notation préfixe / Postfixe. Page Help \\ \hline
Semaine 3 & Vacance (trève de Noel) \\ \hline
Semaine 4 & Vacance (trève de Noel) \\ \hline
Semaine 5 & main et Peaufinage  du projet \\ \hline
\end{tabular}
\end{center}

\chapter*{\vspace*{-2em} Choix de programamtion}
\addcontentsline{toc}{chapter}{Choix de programamtion}

Nous avons décider d'utiliser les symboles + - * pour représenter respectivement l'addition, la soustration et la multiplication. \\ \\
De plus, pour l'option -h (ou -? ou --help), j'ai décidé de réaliser une page man(le manuel sous Linux) en syntaxe \href{https://fr.wikipedia.org/wiki/Roff_(langage_informatique)}{roff} (syntaxe des pages main sous UNIX) pour le coté estétique et pratique. Ce choix induit donc deux dépendances pour le projetqui sont man (pour afficher la syntaxe) et less (dépendance de man). Cependant, il est à noté que ces packages sont installés de base dans la majoritée des systèmes Linux.   

\chapter*{\vspace*{-2em} Bilan personnel pour chaque membre du groupe}
\addcontentsline{toc}{chapter}{Bilan personnel pour chaque membre du groupe}

\section*{Zak Hill}
\addcontentsline{toc}{section}{Zak Hill}
\section*{Serrano Lucas}
\begin{quotation}
"Je trouve personnellement que se projet est trop dirigiste. En effet, devoir compléter des fonctions imposés avec des en-têtes déjà prédéfinis enlève une partie de réflexion sur les choix de programmation. C'est pour une de ces raisons que j'ai décidé de faire une page d'aide en roff (le nom de la syntaxe des pages main), pour retrouvé une certaine liberté créatif."
\emph{— Serrano Lucas}
\end{quotation}
\addcontentsline{toc}{section}{Serrano Lucas}

\end{document}